\newcommand{\TaRcA}{98.4}
\newcommand{\TaRcF}{98.4}
\newcommand{\TaRcP}{0.95}
\newcommand{\TaRgD}{$-1.8$}
\newcommand{\TaRgC}{96.2}
\newcommand{\TaRhD}{$-3.8$}
\newcommand{\TaRhC}{95.4}
\newcommand{\TaRsD}{$-9.6$}
\newcommand{\TaRsC}{89.6}
\newcommand{\TaRsW}{0.68}
\newcommand{\TcRs}{129}
\newcommand{\TcRt}{58}
\newcommand{\TcRo}{138}
\newcommand{\TcRb}{69.0}
\newcommand{\TcRc}{57.5}
\newcommand{\TcRbc}{70.1}
\newcommand{\TaVcA}{98.0}
\newcommand{\TaVcF}{98.0}
\newcommand{\TaVcP}{0.96}
\newcommand{\TaVgD}{$-1.8$}
\newcommand{\TaVgC}{95.8}
\newcommand{\TaVhD}{$-2.0$}
\newcommand{\TaVhC}{96.6}
\newcommand{\TaVsD}{$-6.0$}
\newcommand{\TaVsC}{91.6}
\newcommand{\TaVsW}{0.51}
\newcommand{\TcVs}{191}
\newcommand{\TcVt}{29}
\newcommand{\TcVo}{105}
\newcommand{\TcVb}{86.8}
\newcommand{\TcVc}{67.7}
\newcommand{\TcVbc}{86.8}
\newcommand{\TaCcA}{97.4}
\newcommand{\TaCcF}{97.4}
\newcommand{\TaCcP}{0.14}
\newcommand{\TaCgD}{$-2.8$}
\newcommand{\TaCgC}{95.2}
\newcommand{\TaChD}{$-2.8$}
\newcommand{\TaChC}{94.0}
\newcommand{\TaCsD}{$-20.6$}
\newcommand{\TaCsC}{77.0}
\newcommand{\TaCsW}{0.11}
\newcommand{\TcCs}{220}
\newcommand{\TcCt}{24}
\newcommand{\TcCo}{81}
\newcommand{\TcCb}{90.2}
\newcommand{\TcCc}{75.1}
\newcommand{\TcCbc}{90.1}
\newcommand{\TaZcA}{94.8}
\newcommand{\TaZcF}{94.6}
\newcommand{\TaZcP}{0.93}
\newcommand{\TaZgD}{$-1.8$}
\newcommand{\TaZgC}{94.8}
\newcommand{\TaZhD}{$-1.6$}
\newcommand{\TaZhC}{94.6}
\newcommand{\TaZsD}{$-14.4$}
\newcommand{\TaZsC}{80.2}
\newcommand{\TaZsW}{0.56}
\newcommand{\TcZs}{228}
\newcommand{\TcZt}{17}
\newcommand{\TcZo}{80}
\newcommand{\TcZb}{93.1}
\newcommand{\TcZc}{75.4}
\newcommand{\TcZbc}{93.8}
\newcommand{\TcN}{325}
\newcommand{\TcRej}{0}
\newcommand{\TrepRows}{Grayscale & 0.808 & 0.81 / 0.69 & 96.2 & 0.741 & 0.75 / 0.58 & 95.8 & 0.889 & 0.89 / 0.85 & 95.2 \\
Hue $+108^\circ$ & 0.755 & 0.76 / 0.60 & 95.4 & 0.712 & 0.72 / 0.57 & 96.6 & 0.884 & 0.89 / 0.85 & 94.0 \\
Shift 8\,px & 0.967 & 0.97 / 0.95 & 98.7 & 0.945 & 0.95 / 0.91 & 98.7 & 0.977 & 0.98 / 0.96 & 97.0 \\
Shift 32\,px & 0.929 & 0.93 / 0.91 & 97.9 & 0.943 & 0.94 / 0.87 & 98.7 & 0.965 & 0.97 / 0.94 & 96.9 \\
Patch shuffle & 0.694 & 0.70 / 0.66 & 89.6 & 0.635 & 0.64 / 0.58 & 91.6 & 0.755 & 0.76 / 0.73 & 77.0 \\
Cue conflict & 0.313 & 0.47 / 0.21 & 40.0 & 0.322 & 0.45 / 0.14 & 58.8 & 0.711 & 0.74 / 0.66 & 67.7 \\}

\newcommand{\TOneClean}{All three heads are near ceiling on clean STL-10 (97.4--98.4\% accuracy), and zero-shot CLIP reaches 94.8\%. Later differences are therefore not driven by baseline quality, and we compare both absolute accuracy and change from each model's own clean score. One caveat: the CLIP head's mean confidence of 0.14 is an artefact. It is a linear layer on unit-norm 512-d embeddings with no temperature, so its logits span a small range and its softmax is nearly flat, even though its argmax is accurate. We therefore read CLIP confidence from the zero-shot softmax, which uses the learned logit scale.}
\newcommand{\TOneColor}{Colour is a weak cue for every model. Grayscale costs 1.8--2.8\pp{} and changes at most 5.2\% of predictions. Our hypothesis that hue rotation would hurt more than grayscale held only for ResNet-50 ($-3.8$ vs.\ $-1.8$\pp). ViT and both CLIP predictors lose about the same under hue rotation as under grayscale (within 0.2\pp). The features, however, move a lot (Table~\ref{tab:rep}: $I_T=0.71$--$0.81$ for ResNet/ViT). Chromatic information is still encoded; the heads mostly ignore it.}
\newcommand{\TOneCueSetup}{AdaIN produced 325 candidate conflicts across the ten directions. The count per direction ranges from 22 to 44, because two generation passes wrote overlapping file indices. The rejection rule rejected \emph{no} image: none was a wash-out, and every image kept at least 25\% of its content image's edge energy. It is a weak filter: Fig.~\ref{fig:t1} (first panel) keeps an image whose dog shape is hard to see.}
\newcommand{\TOneCue}{The ranking matches the hypothesis. ResNet-50 is the most texture-sensitive (69.0\% shape bias, 58 texture decisions), ViT-B/16 is intermediate (86.8\%), and CLIP is the most shape-biased (90.2\% for the head, 93.1\% zero-shot). \emph{Coverage limits how far these numbers can be trusted.} ResNet-50 makes a shape-or-texture decision on only 57.5\% of images. Its 138 ``other'' answers go mostly to \emph{bird} (48), \emph{ship} (34) and \emph{deer} (26), classes whose photographs are dominated by background texture (sky, water, foliage). Its bias thus describes only its easiest 57\% of images; CLIP's rests on 75\% coverage and is the best-supported claim. Conditioning on a correctly classified content image changes nothing (last column). The per-pair breakdown is uneven. On airplane--bird, ResNet is texture-dominated (36\% shape bias), and even CLIP drops to 71--78\%. Both classes are photographed against sky, so the style image transfers the other class's \emph{context} along with its texture. Deer--monkey is 90--100\% shape for every model. \emph{Combining colour and cue conflict:} no model relies much on colour, yet ResNet-50 relies on local texture far more than the transformers do. Its texture sensitivity is therefore achromatic (patterns and fur) rather than colour-based.}
\newcommand{\TOneSpatial}{Translation barely matters. With reflection padding, a 32\,px shift (one seventh of the image) costs at most 0.9\pp{} for ResNet-50 and ViT and 2.1--2.4\pp{} for CLIP. Consistency stays at or above 95.6\%, and at 32\,px the four directions differ by at most 1.6\pp. The ViT is as stable as the CNN even though a 32\,px shift moves every token by two patch widths. Random-resized-crop pretraining makes both ImageNet classifiers approximately translation-invariant at the output, whatever the architecture. \emph{Patch shuffling separates the models more:} ResNet-50 loses 9.6\pp, ViT only 6.0\pp, and CLIP the most (20.6\pp{} head, 14.4\pp{} zero-shot). Each $56{\times}56$ tile keeps local parts (a wheel, an eye, fur) that pooled evidence can still vote on; ResNet's average pooling and ViT's attention over permuted tokens both behave like bags of patches, so neither ImageNet classifier needs global arrangement much. Yet confident is not sensible: ResNet-50's errors on shuffled images keep mean confidence 0.68 (27\% above 0.8), committing on the parts and texture that remain; ViT's errors are less confident (0.51).}
\newcommand{\TOneRep}{\emph{(1) Prediction stability overstates representation stability.} Grayscale and hue leave 94--97\% of predictions unchanged, yet ResNet/ViT features move to cosine 0.71--0.81. Nearest-neighbour retrieval shows the information is kept: for 95--97\% of grayscale or hue images, closest clean ResNet/ViT feature is the image's own counterpart. The colour change is encoded, but in directions the head ignores. \emph{(2) The two measures agree on which images move.} Under every intervention, images whose prediction flipped moved further ($I_T^{\neq}<I_T^{=}$). The gap is largest for cue conflict (ResNet 0.47 vs.\ 0.21; ViT 0.45 vs.\ 0.14). \emph{(3) Magnitudes do not compare across backbones.} CLIP has the highest $I_T$ under patch shuffle (0.755), yet shuffling flips 23\% of its predictions. CLIP embeddings lie in a narrow cone, and only 39.6\% of shuffled images retrieve their own clean counterpart (97.4\% for ResNet). The t-SNE (Fig.~\ref{fig:tsne}) shows why: ResNet/ViT shuffled points fall in their clean class clusters, whereas CLIP's form a separate, still class-ordered ``scrambled image'' manifold. Grayscale points sit on their clean clusters for all backbones. \emph{CLIP zero-shot vs.\ its head (RQ3).} The two agree on 96.6\% of clean images, and the head is more accurate (15 vs.\ 2 images that only one of them gets right). Under shuffling, agreement falls to 84\%, and zero-shot wins 40 vs.\ 22. Text prototypes transfer to the displaced manifold better than a head fitted only on clean features. Zero-shot is also the most shape-biased predictor. Same embedding, different decision rule, different cues used.}
\newcommand{\TsynthCues}{In Task~1, texture sensitivity tracked pretraining data and objective more than architecture. Supervised-ImageNet ResNet-50 was most texture-sensitive; the supervised ViT was much less so; and CLIP (a ViT trained with language supervision on web data) was the most shape-biased. Data, augmentation and capacity differ as well, so this is not clean architectural evidence. The same weakness reappears on PACS. Sketch removes colour and texture and keeps contours, and the ERM ResNet-18's largest failures (330 dog$\to$horse, 191 giraffe$\to$horse) are the four-legged animals that share a silhouette and differ mainly in texture and part detail. Classes with distinctive outlines transfer (guitar 93.6\%, house 87.5\%). Colour, which Task~1 showed to be almost unused, is not the problem. Local texture, which ResNet does use, is. This link is qualitative only, since the tasks use different datasets and networks.}
\newcommand{\TtwoStudy}{\emph{Expectation:} larger $\alpha_{\max}$ lowers separability, costs some source accuracy, and helps Sketch up to a point. \emph{Result} (source acc.\ / Sketch acc.\ / Sketch F1 / sep., \%): $\alpha{=}0.25$: 93.1 / 17.1 / 19.5 / 99.9; $\alpha{=}0.5$: 93.2 / 22.7 / 14.1 / 100.0; $\alpha{=}1$: 92.6 / 69.7 / 71.7 / 99.7. Source performance is flat (within 0.6\pp) and separability stays at ceiling, while Sketch accuracy varies by 53\pp. The two weaker settings collapse on Sketch to a single class (\emph{person} at 0.25, 100\% person recall and $\le$10\% elsewhere except guitar; \emph{dog} at 0.5). Only the full reversal kept sketch features inside the region the source head can read. A source-only rule (highest mean source-val F1) would have picked $\alpha{=}0.5$, the second-worst on Sketch, so no source-side signal we logged could have chosen the setting. With one seed per setting, this is evidence of instability, not of an optimal $\alpha$.}
